\chapter{Metodologia ou Material e Métodos}
\addcontentsline{toc}{chapter}{Metodologia ou Material e Métodos}

Este estudo possui uma abordagem qualitativa e quantitativa, e
objetivo exploratório, uma vez que busca investigar a viabilidade
de uma abordagem metodológica. De acordo com \cite{gil2010elaborar}, pesquisas
exploratórias têm como finalidade obter maiores informações sobre o
assunto que se pretende estudar, de maneira a proporcionar uma melhor
delimitação do tema pesquisado, tornando-o mais explícito ou permitindo
a construção de hipóteses. Nesse sentido, o objetivo metodológico não é
estimar um modelo preditivo definitivo, mas avaliar a viabilidade de um procedimento
para transformar comunicações de política monetária em sinais estruturados e
comparáveis ao longo do tempo.

\section{Coleta dos documentos}

O conjunto textual utilizado é composto por documentos oficiais publicados pelo Federal Reserve,
incluindo comunicados das reuniões do Federal Open Market Committee (FOMC), atas,
discursos de dirigentes, testemunhos e comunicados de imprensa. Os documentos foram
coletados a partir do site oficial da instituição. Para cada publicação, foram
armazenadas informações de identificação, como o endereço eletrônico, a data e o
texto, permitindo relacionar os sinais extraídos ao documento de origem e ao período
correspondente.

\section{Extração e classificação dos sinais}

Após a coleta dos documentos, modelos de linguagem de grande porte
(LLMs) foram utilizados para identificar sinais relacionados à política monetária e
organizar as informações extraídas. Cada sinal recebeu uma classificação de postura
monetária, podendo ser hawkish, mostly hawkish, neutral,
mostly dovish ou dovish, além de uma pontuação numérica, uma
justificativa e o trecho de evidência correspondente.

\section{Representação vetorial e construção da rede}

Para representar o conteúdo semântico dos sinais, cada um foi associado a um vetor
de embedding. A partir desses vetores, foi calculada a similaridade do
cosseno entre os sinais. O algoritmo k-nearest neighbors (KNN) foi então
aplicado para selecionar, para cada sinal, os vizinhos semanticamente mais próximos,
respeitando um número máximo de vizinhos e um limiar mínimo de similaridade.

O resultado é uma rede de similaridade semântica ponderada e não direcionada. Seus
nós representam sinais monetários, enquanto suas arestas representam relações
\texttt{SIMILAR\_TO}; o peso de cada aresta corresponde ao escore de similaridade
calculado entre os respectivos vetores. A rede permite identificar proximidades entre
sinais que podem ter sido extraídos de documentos ou datas diferentes, mesmo quando
essa relação não está explicitamente indicada no texto original. As informações de
publicação, data, classificação, pontuação e evidência são mantidas como atributos
dos sinais, permitindo análises posteriores por período.

\section{Análise da rede - REFINAR}

Após a construção da rede, foram aplicados algoritmos de análise de grafos para
investigar sua organização estrutural. O algoritmo de Louvain foi utilizado para
detectar comunidades, isto é, agrupamentos de sinais que apresentam maior
proximidade entre si. A modularidade foi considerada na avaliação da separação entre
essas comunidades e o tamanho dos agrupamentos foi utilizado para caracterizar sua
representatividade na rede.

O PageRank foi empregado como medida de centralidade. Diferentemente de uma simples
contagem de conexões, essa métrica considera também a importância dos nós vizinhos,
permitindo destacar sinais que ocupam posições relevantes na estrutura da rede. As
medidas de centralidade, os agrupamentos identificados e as classificações de postura
monetária foram combinados para produzir um indicador experimental da orientação
presente nas comunicações analisadas.

\section{Comparação temporal e avaliação - REFINAR}

Por fim, o indicador foi organizado por período e comparado com as decisões
observadas do FOMC, especialmente com a direção dos movimentos da taxa de juros. A
avaliação foi realizada por meio da comparação temporal entre a trajetória do
indicador e as decisões efetivamente tomadas, buscando verificar se os sinais
extraídos e sua organização na rede acompanham mudanças na orientação da política
monetária e se apresentam indícios de antecipação dessas mudanças.

Considerando o caráter exploratório da pesquisa, os resultados devem ser
interpretados como evidência inicial sobre a utilidade da abordagem. A análise
privilegia a coerência entre os sinais, a estrutura da rede e os eventos observados,
sem pressupor, nesta etapa, a obtenção de elevado poder preditivo.



% \begin{table}
% \centering
% \caption{Estimates of libero ac erat vulputate feugiat non
%   eu justo. Aenean blandit fermentum efficitur. Fusce faucibus turpis
%   ac dolor sollicitudin, quis lobortis felis porttitor. Phasellus
%   hendrerit quam sit amet erat fringilla, et bibendum ligula feugiat.
%   Phasellus odio tortor, mollis ac cursus ut, fermentum ut orci.
%   Vestibulum malesuada, justo vitae tincidunt posuere, nibh ipsum
%   pellentesque mauris, et suscipit magna turpis pellentesque magna.
%   Mauris sed dui vitae velit facilisis pharetra eu in enim.
%   Suspendisse fringilla lobortis dolor eget sollicitudin.}

% \begin{tabular}{ c p{5.5pc} c p{5.5pc} }
% \hline
% \multicolumn{2}{c}{Múltiplas colunas} & Aqui & Acolá \\
% \hline
% Um & \raggedright\arraybackslash Largura fixa em (5.5pc). & Três &
% \raggedright\arraybackslash Coluna 4, largura fixa.\\
% \hline
% \label{Tab:estim}
% \end{tabular}
% \end{table}

